\section{Resolution in Plan Manifests}
\label{sec:manifests}

\subsection{Motivation}

The \redist{} audit chain~\citep{dellaLibera2026b1} records a cryptographic
SHA-256 hash of every plan output together with the parameters used to
generate it.
This chain enables independent verification: given the adjacency file,
the YAML config, and the manifest, a verifier can reproduce the run and
confirm the hash.

Resolution introduces two new verification requirements:
\begin{enumerate}
  \item A verifier must know at which geographic level the plan was generated
        (tract, BG, or county) in order to interpret the assignment file
        correctly.
  \item For multi-scale runs, a verifier must be able to reconstruct the
        fine-to-coarse partition from public data, without running the
        pipeline code.
\end{enumerate}

Both requirements are met by three new fields added to every plan manifest
(\texttt{index.json}) and four additional fields for multi-scale runs.

\subsection{Core Resolution Fields}

Every plan manifest gains the following fields at generation time:

\begin{lstlisting}[language=json]
"plan_resolution": "tract",
"n_units": 2195,
"unit_type": "census tract"
\end{lstlisting}

\noindent
\texttt{plan\_resolution} takes one of three values: \texttt{"bg"},
\texttt{"tract"}, or \texttt{"county"}.
\texttt{n\_units} is the number of geographic units at the plan resolution
(number of rows in the assignment file).
\texttt{unit\_type} is a human-readable label used in reports and analysis
output.

These fields are written by \texttt{run\_single\_state()} at plan generation
time and read by \texttt{redist analyze} and \texttt{redist report} to
interpret the assignment file correctly.
For single-scale runs (the current default), \texttt{plan\_resolution = "tract"}
and \texttt{n\_units} equals the number of tracts in the state.

\subsection{Multi-Scale Manifest Fields}

For multi-scale runs, four additional fields are written:

\begin{lstlisting}[language=json]
"multiscale_fine": "tract",
"multiscale_coarse": "county",
"fine_to_coarse_formula": "geoid_prefix[:5]",
"fine_to_coarse_comment":
  "first 5 chars of 11-char tract GEOID = parent county FIPS"
\end{lstlisting}

\noindent
The \texttt{fine\_to\_coarse\_formula} field encodes the partition derivation
rule as a short string that a verifier can apply manually.
Combined with the \texttt{index\_to\_geoid\_file} field (present in all
manifests, pointing to the GEOID lookup file co-located with the adjacency
file), a verifier can:
\begin{enumerate}
  \item Load the index-to-GEOID map from \texttt{index\_to\_geoid\_file}.
  \item Apply the prefix rule: for each unit index $i$, take the prefix
        of length specified by \texttt{fine\_to\_coarse\_formula}.
  \item Look up the prefix in the coarse-level GEOID set to obtain the
        coarse index.
\end{enumerate}
This reconstruction is possible without any pipeline code --- only the
public adjacency files and a simple string-prefix operation are required.

\subsection{Manifest Examples for All Three Options}

\subsubsection{Option A: BG$\to$Tract}

\begin{lstlisting}[language=json]
{
  "plan_resolution": "bg",
  "n_units": 5012,
  "unit_type": "census block group",
  "index_to_geoid_file": "nc_bg_adjacency_2020_geoids.json",
  "multiscale_fine": "bg",
  "multiscale_coarse": "tract",
  "fine_to_coarse_formula": "geoid_prefix[:11]",
  "fine_to_coarse_comment":
    "first 11 chars of 12-char BG GEOID = parent tract GEOID"
}
\end{lstlisting}

\subsubsection{Option B: Tract$\to$County}

\begin{lstlisting}[language=json]
{
  "plan_resolution": "tract",
  "n_units": 2195,
  "unit_type": "census tract",
  "index_to_geoid_file": "nc_adjacency_2020_geoids.json",
  "multiscale_fine": "tract",
  "multiscale_coarse": "county",
  "fine_to_coarse_formula": "geoid_prefix[:5]",
  "fine_to_coarse_comment":
    "first 5 chars of 11-char tract GEOID = parent county FIPS"
}
\end{lstlisting}

\subsubsection{Option C: BG$\to$County}

\begin{lstlisting}[language=json]
{
  "plan_resolution": "bg",
  "n_units": 5012,
  "unit_type": "census block group",
  "index_to_geoid_file": "nc_bg_adjacency_2020_geoids.json",
  "multiscale_fine": "bg",
  "multiscale_coarse": "county",
  "fine_to_coarse_formula": "geoid_prefix[:5]",
  "fine_to_coarse_comment":
    "first 5 chars of 12-char BG GEOID = parent county FIPS"
}
\end{lstlisting}

\noindent
Note that Options B and C share the same prefix length (5) but differ in
the source GEOID length (11 for tracts, 12 for BGs) and in
\texttt{plan\_resolution}.

\subsection{Non-Multi-Scale Runs}

For single-resolution runs, the fields \texttt{multiscale\_fine},
\texttt{multiscale\_coarse}, \texttt{fine\_to\_coarse\_formula},
and \texttt{fine\_to\_coarse\_comment} are absent from the manifest.
The core resolution fields (\texttt{plan\_resolution}, \texttt{n\_units},
\texttt{unit\_type}) are always present regardless of whether multi-scale
was used.

\subsection{Downstream Pipeline Behavior}

\textbf{Analysis} (\texttt{redist label-analyze}): reads
\texttt{plan\_resolution} from the manifest to determine which population
file and geometry file to use for compactness computation.
BG-level plans use BG geometries for Polsby-Popper~\citep{polsby1991third};
tract-level plans use tract geometries (current default).
County-level plans would use county geometries (deferred).

\textbf{Reporting} (\texttt{redist label-report}): uses \texttt{unit\_type}
from the manifest for the human-readable label in HTML and JSON reports
(e.g., ``district boundary at census tract resolution'' vs.\ ``\ldots at
census block group resolution'').

\textbf{Verification} (\texttt{redist label-verify}): the SHA-256 hash
covers the entire manifest including resolution fields.
A plan rebuilt at a different resolution will not match the stored hash,
flagging a resolution mismatch as a verification failure.
